% !TeX root = ../papers/nc-p3.tex
\section{種数2の四窓表示と全ての有理縦切片}
\label{sec:rank-fifteen-genus-two-unit-line}

\begin{proposition}\label{prop:rank-fifteen-diagonal-genus-two}
標数3の代数閉体 $k$ 上で $X\ne0,1$ を固定し、前節の $b=1/c$, $D=v^3-v$,
$d=b^3-b$ を使う。幾何学的一般点で四窓は種数2の体
$k(v,\sqrt{D^2-d^2})$ に係数を持つKlein斉次な表示と添字保持同型である。
この表示は $D(D^2-d^2)\ne0$ の有限点で使える。
\end{proposition}
\begin{proof}
前節の四位置のWalsh変換を使い、対角倍率はまだ掛けない。
$a=1$, $b_0=(K-1)/2$ とし、六根の組を
$(p_0,m_0)=(L+M,L-M)$, $(p_1,m_1)=(B-A,B+A)$,
$(p_2,m_2)=(D_0-C,D_0+C)$ とする。ここで $D_0$ は正規形の最後の倍率である。
共通定数を除く非零成分は36個で、各成分 $u_w$ は符号を除き、
$a,b_0,a+b_0$ と六根のうちの一つである。各語 $w$ の四位置の矢を示す
$0,1$ 行列を $M$ とし、辞書順に最初の独立な13行を $I$ とする。
$M_I$ の最初の独立な13列の小行列式は $-4$ で、$M=N M_I$ の係数 $N$ は全て整数である。
従って代数閉包上の対角基底変更で $I$ の13係数を同時に1とし、他の係数を
$u_w\prod_{r\in I}u_r^{-N_{wr}}$ とできる。実際、対応するトーラス準同型は
階数13で全射である。この操作は四位置で周期的である。

\path{computations/verify_rank_fifteen_genus_two_unit_line.py} はこの行列を四窓の式から生成する。
規格化した各単項式では六根の指数が全て偶数か全て奇数になる。
例えば語 $0231$ の係数は $-m_0m_1p_2/(m_2p_0p_1)$ である。
従って必要なのは六根の二乗とその積 $W$ のみである。
前節の三次式の積より $W^2$ の平方類は $D^2-d^2$ である。
これは相異なる六単根を持つので、二重被覆の種数は2である。
除算に現れる元は六根および $a,b_0,a+b_0$ に限る。
後二者の零では $K=\pm1$、従って $D=0$ であり、主張した開集合で変換は可逆である。
\end{proof}

\begin{theorem}\label{thm:rank-fifteen-all-negative-unit-residues}
剰余標数3の付値で $\bar x=-1$ なら、任意の非零の $y$ で通常な標数零の階数15窓は
標準Koszul左右AS正則で、各始点は四中心四次式上自由階数256である。
\end{theorem}
\begin{proof}
正付値・負付値の $y$ は既知なので単元を扱う。$i^2=-1$, $K=1/v$,
$H=i(K-v)$ とすると $Y=\bar y=-D^2$ である。上の六根の二乗は順に
\[
 (K^2-1\pm2H)/4,\qquad (K+1\pm H)/2,\qquad (K-1\pm H)/2 .
\]
$Z^2=D^2+1$ と置けば $W=Z D^2/v^6$ となり、二乗して六式の積と照合できる。
従って $\mathbb F_3(v)$ 上 $1,i,Z,iZ$ の四座標で演算できる。
符号に応じた $Z$ の選択は全て含まれる。

上のコードで五次まで経路基底を作り、変換後の文字別の四次中心次元が
$(5,0,0,0)$ になることを関数体上で求める。文字0の自由列の昇順の五本のうち
後の四本を課す。全始点で十三次まで中心商の次元は $(1+t+t^2+t^3)^4$ の係数に一致する。
各逆元の分子の $\mathbb F_3(v)$ へのノルムの零点を記録する。
入力・経路・中心を含めた各始点の平方自由な除外多項式の次数は
$1049,2866,1642,3847$、最小公倍式 $F$ の次数は8801である。

$Y=0,1$ を除くと各水準の六表示は $v\mapsto\pm v+j$ ($j\in\mathbb F_3$) の一軌道である。
$F$ からこの二水準の因子を除き、六変換による引き戻しの最大公約元を取ると次数156となる。
各既約因子で $Y$ の最小多項式 $g$ を求め、同じ $g$ の因子の積が
$g(-D^2)$ のモニック化に等しいことを確認する。全表示が除かれる八軌道の次数は
$2,2,2,3,3,4,4,6$ である。四次までと $Y=1$ は既知で、残る六次式は
\[
 Y^6+2Y^5+Y^4+2Y^3+Y+1 .
\]
この根に対し $\mathbb F_{3^{12}}$ で元の六根の表示を取り、四始点で長さ13まで直接検算する。
Frobenius共役でその六根を覆う。超越的な $Y$ はどの除外多項式にも当たらない。

選んだ剰余表示は定理\ref{thm:rank-fifteen-unit-residues}のモニックな方程式で持ち上がる。
上の基底変更の可逆行列も単元行列式のまま持ち上げる。
経路の一般下界と四次中心の全次元の下界5により、剰余の四つの中心式が持ち上がる。
十三次消滅を含む中心商の上界と既出の自由性判定から全次数の結論が従う。
\end{proof}

\begin{corollary}
標数零の代数閉体で、任意の有理数 $x$ と任意の非零の幾何水準 $y$ において、
通常な階数15窓は全て上の正則性と自由性を持つ。
さらに3の上のある付値で $\bar a=-1$ となる縦切片 $x=a$ では全非零水準の
添字保持同型類は系\ref{cor:rank-fifteen-unit-vertical-conics}の $T$ の式で一意に表される。
$a^3/27,(1-a)^2/4$ は各二類、その他は各四類で、最小周期は全て2である。
\end{corollary}
\begin{proof}
有理数 $x$ は3で負付値を持つか、剰余が $0,1,-1$ である。付値を係数体に延長し、
既知の三領域と本節の定理を適用する。$\bar a=-1$ では $a=3,3/4$ および
$4a^2-12a+7=0$ の二根はいずれも生じない。従って通常表示が全点を覆い、
二分岐水準は異なる。円錐曲線の表示と最小周期は前掲の系と同じ計算で得られる。
\end{proof}
